\chapter{Nephritic Syndrome}
\index{Nephritic Syndrome}
\label{sec:Nephritic Syndrome}

\section{Overview}
Nephritic syndrome is a clinical complex resulting from acute glomerular inflammation characterized by leukocyte infiltration, cellular growth, and reactive glomerular capillary wall injury. This genetic disorder results from specific gene mutations or chromosomal abnormalities. It may be present at birth or manifest later in life depending on the underlying mechanism. Genetic disorders result from abnormalities in an individual's DNA, ranging from single-gene mutations to chromosomal abnormalities. These conditions may be inherited from parents or arise from new mutations. Pattern of inheritance depends on whether the mutation is dominant, recessive, or X-linked. Genetic counseling helps families understand risks and make informed decisions.
\section{Causes}
This condition happens because of immune complex deposition or antibody binding within the glomerulus, causing a decrease in glomerular filtration rate (GFR) and breakdown of the capillary barrier to red blood cells and protein. Etiologies include poststreptococcal glomerulonephritis (PSGN), IgA nephropathy (Berger disease), lupus nephritis, and anti-GBM disease (Goodpasture syndrome). Genetic disorders result from mutations in specific genes that encode proteins essential for normal cellular function. The pattern of inheritance depends on whether the mutation is dominant, recessive, or X-linked.   
\section{Risk Factors}

\begin{itemize}[noitemsep]
  \item Genetic predisposition and family history
  \item Advancing age
  \item Lifestyle factors (diet, exercise, smoking, alcohol)
  \item Environmental exposures
  \item Underlying medical conditions
  \item Medication use
\end{itemize}
\section{Signs and Symptoms}

\subsection{Common symptoms}
\begin{itemize}[noitemsep]
  \item You may experience the classic triad of hematuria ('cola-colored' or dysmorphic RBCs and RBC casts), high blood pressure (due to fluid retention), and oliguria with mild-to-moderate proteinuria ($<3.5$ g/day) and facial/periorbital swelling. Diagnosis is supported by urinalysis showing hematuria and proteinuria, elevated serum creatinine, low serum complement levels (in PSGN and lupus), and definitive kidney biopsy.

  \item Pain or discomfort in the affected area
  \end{itemize}

\subsection{Emergency warning signs}
\begin{itemize}[noitemsep]
  \item Severe or worsening symptoms of nephritic syndrome require immediate medical evaluation
\end{itemize}
\section{Progression and Stages}
The condition may progress through different stages of severity over time. Early stages often involve mild or subtle symptoms that may be overlooked. Moderate stages involve more noticeable symptoms affecting daily function. Advanced stages may involve significant impairment and complications.
\section{Complications}
\begin{itemize}[noitemsep]
  \item Disease progression leading to worsening symptoms
  \item Secondary infections or injuries
  \item Side effects of treatment
  \item Reduced quality of life
  \item Emotional and psychological impact
\end{itemize}
\section{Diagnosis}
\begin{itemize}[noitemsep]
  \item Comprehensive medical history and physical examination
  \item Laboratory tests to assess organ function
  \item Imaging studies to visualize affected structures
  \item Specialized testing depending on the affected system
  \item Consultation with relevant specialists
\end{itemize}
\section{Prevention}
\begin{itemize}[noitemsep]
  \item Maintain a healthy lifestyle with balanced diet and exercise
  \item Avoid known risk factors
  \item Attend regular medical check-ups for early detection
  \item Follow recommended screening guidelines
\end{itemize}
\section{Treatment Options}
Treatment focuses on fluid and sodium restriction, antihypertensive therapy (diuretics, ACE inhibitors), and targeted immunosuppressants (corticosteroids, cyclophosphamide) for rapidly progressive cases. Symptomatic management and supportive care Enzyme replacement therapy for certain conditions Gene therapy (emerging for select conditions) Dietary modifications (e.g., phenylketonuria) Regular monitoring for associated complications Physical and occupational therapy Genetic counseling for family planning Participation in clinical trials for new therapies

\begin{itemize}[noitemsep]
  \item Medications to manage symptoms and address underlying mechanisms
  \item Lifestyle modifications and self-care measures
  \item Physical therapy and rehabilitation
  \item Surgical intervention when indicated
  \item Regular monitoring and follow-up care
\end{itemize}
\section{Living with It}
\begin{itemize}[noitemsep]
  \item Follow your treatment plan as prescribed
  \item Maintain regular follow-up with your healthcare provider
  \item Make necessary lifestyle adjustments
  \item Seek support from family, friends, or support groups
  \item Stay informed about your condition
\end{itemize}
\section{Outlook}
Your outlook depends on underlying histology.
